\documentclass[12pt,a4paper]{article}
\usepackage[utf8]{inputenc}
\usepackage[italian]{babel}
\usepackage{geometry}
\usepackage{graphicx}
\usepackage{listings}
\usepackage{xcolor}
\usepackage{hyperref}
\usepackage{booktabs}
\usepackage{array}
\usepackage{amsmath}

\geometry{margin=2.5cm}

\definecolor{codegreen}{rgb}{0,0.6,0}
\definecolor{codegray}{rgb}{0.5,0.5,0.5}
\definecolor{codepurple}{rgb}{0.58,0,0.82}
\definecolor{backcolour}{rgb}{0.95,0.95,0.92}

\lstdefinestyle{javastyle}{
    backgroundcolor=\color{backcolour},
    commentstyle=\color{codegreen},
    keywordstyle=\color{blue},
    numberstyle=\tiny\color{codegray},
    stringstyle=\color{codepurple},
    basicstyle=\ttfamily\footnotesize,
    breakatwhitespace=false,
    breaklines=true,
    captionpos=b,
    keepspaces=true,
    numbers=left,
    numbersep=5pt,
    showspaces=false,
    showstringspaces=false,
    showtabs=false,
    tabsize=2,
    language=Java
}

\title{\textbf{Analisi della Scelta Progettuale: Pattern Mapper}\\[0.5cm]
\large Applicazione dei Principi GRASP nel Design delle Classi di Conversione}
\author{Progetto ISPW - Habibi Yalla}
\date{Febbraio 2026}

\begin{document}

\maketitle

\tableofcontents
\newpage

%==============================================================================
\section{Introduzione}

Nel contesto dello sviluppo del sistema \textbf{Habibi Yalla}, si è resa necessaria l'introduzione di classi dedicate alla conversione tra Entity e Bean: i \textbf{Mapper}. Questo documento analizza la scelta progettuale dal punto di vista dei principi \textbf{GRASP} (General Responsibility Assignment Software Patterns), evidenziando benefici e trade-off architetturali.

\subsection{Contesto Architetturale}

Il sistema adotta l'architettura \textbf{BCE} (Boundary-Control-Entity) raffinata in \textbf{MVC}. In questa architettura:
\begin{itemize}
    \item Le \textbf{Entity} rappresentano il dominio di business (es. \texttt{Food}, \texttt{Voucher}, \texttt{Ordine})
    \item I \textbf{Bean} sono oggetti di trasporto dati (DTO) tra i layer Boundary e Control
    \item La conversione Entity $\leftrightarrow$ Bean è necessaria per il disaccoppiamento dei layer
\end{itemize}

%==============================================================================
\section{Problema: Dove Posizionare la Logica di Conversione?}

\subsection{Situazione Iniziale}

Prima dell'introduzione dei Mapper, i metodi di conversione erano posizionati all'interno dei \textbf{Controller Applicativi}:

\begin{lstlisting}[style=javastyle, caption={Conversione nel Controller (PRIMA)}]
public class CreaOrdineController {
    
    // Responsabilita 1: Business Logic
    public OrdineBean inizializzaNuovoOrdine(String clienteId) { ... }
    public boolean aggiungiProdottoAOrdine(FoodBean foodBean) { ... }
    
    // Responsabilita 2: Conversione (INCOERENTE!)
    private FoodBean convertFoodToBean(Food food) {
        FoodBean bean = new FoodBean();
        bean.setId(food.getId());
        bean.setDescrizione(food.getDescrizione());
        // ... altri campi
        return bean;
    }
    
    private List<FoodBean> convertFoodListToBeanList(List<Food> foods) {
        // ... iterazione e conversione
    }
}
\end{lstlisting}

\subsection{Problemi Identificati}

\begin{enumerate}
    \item \textbf{Violazione High Cohesion}: Il Controller aveva due responsabilità distinte
    \item \textbf{Codice Duplicato}: La stessa logica di conversione era presente in più Controller
    \item \textbf{Difficoltà di Manutenzione}: Modifiche alla struttura del Bean richiedevano interventi in più punti
\end{enumerate}

%==============================================================================
\section{Soluzione: Introduzione dei Mapper}

\subsection{Design della Soluzione}

Sono state create tre classi Mapper dedicate:

\begin{table}[h]
\centering
\begin{tabular}{|l|l|l|}
\hline
\textbf{Mapper} & \textbf{Entity} & \textbf{Bean} \\
\hline
\texttt{FoodMapper} & \texttt{Food} & \texttt{FoodBean} \\
\texttt{VoucherMapper} & \texttt{Voucher} & \texttt{VoucherBean} \\
\texttt{OrdineMapper} & \texttt{Ordine} & \texttt{OrdineBean} \\
\hline
\end{tabular}
\caption{Mapping Entity-Bean}
\end{table}

\subsection{Struttura del FoodMapper}

\begin{lstlisting}[style=javastyle, caption={FoodMapper - Implementazione}]
public class FoodMapper {

    private FoodMapper() {
        // Utility class - previene istanziazione
    }

    public static FoodBean toBean(Food food) {
        if (food == null) {
            return null;
        }
        FoodBean bean = new FoodBean();
        bean.setId(food.getId());
        bean.setDescrizione(food.getDescrizione());
        bean.setCosto(food.getCosto());
        bean.setDurata(food.getDurata());
        bean.setTipo(food.getTipo());
        bean.setClasse(food.getClass().getSimpleName());
        return bean;
    }

    public static List<FoodBean> toBeanList(List<Food> foodList) {
        if (foodList == null) {
            return new ArrayList<>();
        }
        List<FoodBean> beans = new ArrayList<>();
        for (Food food : foodList) {
            FoodBean bean = toBean(food);
            if (bean != null) {
                beans.add(bean);
            }
        }
        return beans;
    }
}
\end{lstlisting}

%==============================================================================
\section{Analisi GRASP}

\subsection{Pure Fabrication}

Il pattern \textbf{Pure Fabrication} è il principio GRASP primario applicato ai Mapper.

\begin{quote}
\textit{``Una Pure Fabrication è una classe che non rappresenta un concetto del dominio di business, ma viene creata per ottenere alta coesione, basso accoppiamento e riuso.''}
\end{quote}

\subsubsection{Applicazione ai Mapper}

\begin{itemize}
    \item \textbf{Non è un concetto di dominio}: ``FoodMapper'' non esiste nel vocabolario del business (kebab, ordini, voucher)
    \item \textbf{Creata artificialmente}: Introdotta esclusivamente per migliorare la qualità del design
    \item \textbf{Responsabilità tecnica}: Gestisce solo la trasformazione strutturale dei dati
\end{itemize}

\subsubsection{Vantaggi della Pure Fabrication}

\begin{enumerate}
    \item Evita di ``inquinare'' le Entity con logica di presentazione
    \item Mantiene i Controller focalizzati sulla business logic
    \item Crea un punto centralizzato per la gestione delle conversioni
\end{enumerate}

\subsection{High Cohesion}

\subsubsection{Definizione}

\begin{quote}
\textit{``Una classe ha alta coesione quando tutte le sue responsabilità sono fortemente correlate e focalizzate su un singolo scopo.''}
\end{quote}

\subsubsection{Impatto sui Controller}

\begin{table}[h]
\centering
\begin{tabular}{|l|c|c|}
\hline
\textbf{Metrica} & \textbf{Prima} & \textbf{Dopo} \\
\hline
Responsabilità CreaOrdineController & 2 & 1 \\
Metodi di conversione nel Controller & 2 & 0 \\
Linee di codice Controller & 262 & 234 \\
Focus del Controller & Misto & Solo Business Logic \\
\hline
\end{tabular}
\caption{Impatto sui Controller Applicativi}
\end{table}

\subsubsection{Coesione dei Mapper}

I Mapper hanno \textbf{coesione molto alta} perché:
\begin{itemize}
    \item Tutti i metodi riguardano la conversione di un'unica coppia Entity-Bean
    \item Non hanno dipendenze esterne (solo i tipi da convertire)
    \item Non mantengono stato (metodi statici, utility class)
\end{itemize}

\subsection{Low Coupling}

\subsubsection{Definizione}

\begin{quote}
\textit{``Una classe ha basso accoppiamento quando dipende da poche altre classi e le dipendenze sono stabili.''}
\end{quote}

\subsubsection{Analisi delle Dipendenze}

\textbf{Prima (Controller con conversione interna):}
\begin{itemize}
    \item Controller $\rightarrow$ Entity (Food)
    \item Controller $\rightarrow$ Bean (FoodBean)
    \item Controller $\rightarrow$ Facade, DAO, ecc.
\end{itemize}

\textbf{Dopo (con Mapper):}
\begin{itemize}
    \item Controller $\rightarrow$ Mapper (dipendenza stabile)
    \item Controller $\rightarrow$ Bean (FoodBean)
    \item Mapper $\rightarrow$ Entity + Bean
\end{itemize}

\subsubsection{Valutazione}

Il Mapper introduce una nuova dipendenza, ma questa è \textbf{accettabile} perché:
\begin{enumerate}
    \item Il Mapper è una classe \textbf{stabile} (cambia raramente)
    \item La dipendenza è \textbf{unidirezionale} (Controller $\rightarrow$ Mapper)
    \item Riduce la \textbf{conoscenza} che il Controller deve avere dell'Entity
\end{enumerate}

\subsection{Information Expert}

\subsubsection{Il Trade-off}

Il principio Information Expert afferma:
\begin{quote}
\textit{``Assegna la responsabilità alla classe che ha le informazioni necessarie per adempierla.''}
\end{quote}

Applicato rigorosamente, la conversione dovrebbe essere nell'Entity:

\begin{lstlisting}[style=javastyle, caption={Alternativa: Conversione nell'Entity}]
public class Food {
    // ... attributi ...
    
    public FoodBean toBean() {
        FoodBean bean = new FoodBean();
        bean.setId(this.id);
        bean.setDescrizione(this.descrizione);
        // ... l'Entity conosce i propri dati
        return bean;
    }
}
\end{lstlisting}

\subsubsection{Perché Non Usare Information Expert}

Questa soluzione \textbf{viola la separazione BCE}:
\begin{itemize}
    \item L'Entity (layer Model) conoscerrebbe il Bean (layer Boundary/Control)
    \item Creerebbe un \textbf{coupling architetturale} tra layer che dovrebbero essere indipendenti
    \item Renderebbe l'Entity dipendente dalla struttura di presentazione
\end{itemize}

\subsubsection{Conclusione}

\textbf{Pure Fabrication prevale su Information Expert} in questo caso, per preservare la stratificazione architetturale BCE/MVC.

\subsection{Controller (GRASP)}

Il pattern GRASP Controller stabilisce che:
\begin{quote}
\textit{``Il Controller coordina e delega, non dovrebbe contenere logica di business complessa o responsabilità secondarie.''}
\end{quote}

\subsubsection{Prima del Refactoring}

\begin{lstlisting}[style=javastyle]
public List<FoodBean> getProdottiBaseDisponibili() {
    List<Food> foodBase = FoodLazyFactory.getInstance().getAllFoodBase();
    return convertFoodListToBeanList(foodBase); // Conversione INTERNA
}
\end{lstlisting}

\subsubsection{Dopo il Refactoring}

\begin{lstlisting}[style=javastyle]
public List<FoodBean> getProdottiBaseDisponibili() {
    List<Food> foodBase = FoodLazyFactory.getInstance().getAllFoodBase();
    return FoodMapper.toBeanList(foodBase); // DELEGA al Mapper
}
\end{lstlisting}

Il Controller ora \textbf{delega} invece di implementare direttamente.

%==============================================================================
\section{Riepilogo dei Principi GRASP}

\begin{table}[h]
\centering
\begin{tabular}{|l|c|l|}
\hline
\textbf{Principio} & \textbf{Applicato?} & \textbf{Note} \\
\hline
Pure Fabrication & \checkmark & Principio primario dei Mapper \\
High Cohesion & \checkmark & Controller e Mapper più coesi \\
Low Coupling & \checkmark & Dipendenze stabili \\
Information Expert & $\times$ & Violato per preservare BCE \\
Controller & \checkmark & Delega invece di implementare \\
\hline
\end{tabular}
\caption{Riepilogo applicazione principi GRASP}
\end{table}

%==============================================================================
\section{Conclusioni}

L'introduzione delle classi Mapper rappresenta una scelta progettuale conforme ai principi GRASP, in particolare:

\begin{enumerate}
    \item \textbf{Pure Fabrication} giustifica la creazione di classi non legate al dominio
    \item \textbf{High Cohesion} migliora sia nei Controller che nei nuovi Mapper
    \item \textbf{Low Coupling} è mantenuto grazie a dipendenze stabili e unidirezionali
    \item Il \textbf{trade-off} con Information Expert è accettabile per preservare l'architettura BCE
\end{enumerate}

Il risultato è un design più manutenibile, testabile e conforme ai principi di buona progettazione object-oriented.

\end{document}
